\documentclass{article}
\usepackage{import}
\import{../../template/}{article-template}
\graphicspath{{../../images/}}
\addbibresource{../../template/library.bib}

\title{Javascript}
\author{PENG}
\date{\today}

\begin{document}
\maketitle
\tableofcontents
\section{Introduction}
\url{https://developer.mozilla.org/en-US/docs/Web/JavaScript/Guide/}

\begin{definition}[javascript]
	Javascript is an object-oriented and dynamically typed scripting language used to make webpages interactive.
\end{definition}

\begin{remark}~
	\begin{itemize}
		\item Always write a semicolon after a statement.
		\item Variables should always be declared before they are used.
	\end{itemize}
\end{remark}

\begin{definition}[declaration]~
	\begin{itemize}
		\item \verb|var|
		\item \verb|let| : block-scoped
		\item \verb|const| : block-scoped, must be initialized
	\end{itemize}
\end{definition}

\begin{example}[var scope]~
	\begin{verbatim}
console.log(x === undefined); // true
var x = 3;
(function () {
    console.log(x); // undefined
    var x = "local value";
})();	
\end{verbatim}
\end{example}

\begin{definition}[types]~
	\begin{itemize}
		\item \verb|Boolean| : true or false
		\item \verb|null|
		\item \verb|Number| : integer or floating point number
		\item \verb|String|
	\end{itemize}
\end{definition}

\begin{definition}[type conversion]~
	\begin{itemize}
		\item num $\to$ string : \verb|"a is " + 10 // "a is 10"|
		\item string $\to$ num : \verb|parseInt("10.1")|
	\end{itemize}
\end{definition}

\printbibliography
\end{document}
